\section{Two-Stage KAN Controller}\label{sec:controller}

\subsection{Variable-$N$ Architecture}

The cluster controller must predict $(\mu, \beta_A, \chi)$ at the cluster level and $\{\beta_i\}_{i=1}^N$ at the user level, where $N$ may vary across deployments. We adopt a two-stage architecture (Fig.~\ref{fig:ctrl}) that handles arbitrary $N$ by parameter sharing:

\begin{itemize}
\item \emph{Global head} $\mathcal{H}_g$: input $\bs_g = (s_A, N, \bar s_B, \min_i s_i, d_E)$; output $(\hat\mu, \hat\chi, \hat\beta_A)$.
\item \emph{Per-user shared head} $\mathcal{H}_u$: input $(s_i, s_A, \hat\mu, \hat\chi)$; output $\hat\beta_i$. The same head is invoked once per user $i$.
\end{itemize}

\begin{figure}[!htbp]
\centering
\resizebox{\columnwidth}{!}{% Two-stage controller (Fig. 3): global head -> per-user shared head.
%
% Clean left-to-right dataflow:
%   cluster-aggregate features -> global head H_g -> (mu,chi,beta_A)
%   then a SHARED per-user head H_u taking (s_i, s_A, mu_hat, chi_hat)
%   fanning out to users B_1..B_N to make parameter-sharing-over-N obvious.
% Muted palette: inputs black!6, global head blue!8, per-user head teal!14,
% outputs orange!16; shared head wrapped in a light backgrounds rectangle.
\begin{tikzpicture}[
  font=\footnotesize,
  >={Latex[length=1.6mm,width=1.1mm]},
  ablk/.style  ={draw=black!55,thick,fill=black!6,text=black,
                 rounded corners=2pt,align=center,inner sep=4pt},
  ghead/.style ={draw=black!55,thick,fill=blue!8,text=black,
                 rounded corners=2pt,align=center,inner sep=5pt,
                 minimum height=12mm},
  uhead/.style ={draw=black!55,thick,fill=teal!14,text=black,
                 rounded corners=2pt,align=center,inner sep=5pt,
                 minimum height=12mm},
  obox/.style  ={draw=black!55,thick,fill=orange!16,text=black,
                 rounded corners=2pt,align=center,inner sep=4pt},
  ubox/.style  ={draw=black!55,thick,fill=orange!16,text=black,
                 rounded corners=2pt,align=center,inner sep=3pt,
                 minimum width=11mm,minimum height=6.5mm},
  ghost/.style ={draw=black!30,thick,fill=teal!7,text=black!55,
                 rounded corners=2pt,align=center,inner sep=4pt,
                 minimum width=11mm,minimum height=6.5mm},
  flow/.style  ={->,semithick,draw=black!70},
  fan/.style   ={->,semithick,draw=orange!75!black},
  glab/.style  ={font=\scriptsize,text=black!65,inner sep=1.5pt},
]

% ===================== Stage 1: global head ==============================
\node[ablk] (FG) at (0,0)
  {Cluster-aggregate\\features\\[1pt]
   {\scriptsize$\bs_g=(s_A,N,\bar s_B,\min_i s_i,d_E)$}};

\node[ghead,right=14mm of FG] (HG)
  {\textbf{Global head}~$\mathcal{H}_g$\\[1pt]{\scriptsize KAN / MLP / Linear}};

\node[obox,right=14mm of HG] (OG)
  {$(\hat\mu,\;\hat\chi,\;\hat\beta_A)$};

\draw[flow] (FG) -- (HG);
\draw[flow] (HG) -- (OG);

% ===================== Stage 2: shared per-user head =====================
% Per-user feature input, sitting below the global row.
\node[ablk,below=20mm of FG] (FU)
  {Single-user\\features\\[1pt]{\scriptsize$(s_i,\,s_A)$}};

\node[uhead,below=20mm of HG] (HU)
  {\textbf{Per-user head}~$\mathcal{H}_u$\\[1pt]
   {\scriptsize shared weights}\\[1pt]
   {\scriptsize input $(s_i,s_A,\hat\mu,\hat\chi)$}};

\draw[flow] (FU) -- (HU);

% Predicted global block (mu_hat, chi_hat) feeds the per-user head from above.
% Route: down out of OG, left to HU's x, down into HU.north (strict orthogonal).
\coordinate (Gd) at ($(OG.south)+(0,-6mm)$);
\draw[flow] (OG.south) -- (Gd) -| (HU.north)
  node[glab,pos=0.32,above]{$\hat\mu,\hat\chi$ (global block)};

% ----- Fan-out to per-user outputs B_1..B_N (parameter sharing) ----------
\node[ubox,right=16mm of HU]            (B1) {$\hat\beta_1$};
\node[ubox,below=3.5mm of B1]           (B2) {$\hat\beta_2$};
\node[ghost,below=3.5mm of B2]          (Bd) {$\vdots$};
\node[ubox,below=3.5mm of Bd]           (BN) {$\hat\beta_N$};

\draw[fan] (HU.east) -- (B1.west);
\draw[fan] (HU.east) -- (B2.west);
\draw[fan] (HU.east) -- (Bd.west);
\draw[fan] (HU.east) -- (BN.west);

% Brace marking one shared head applied N times.
\draw[decorate,decoration={brace,amplitude=4pt,mirror},draw=black!55]
  ($(B1.north east)+(3mm,0)$) -- ($(BN.south east)+(3mm,0)$)
  node[midway,right=5pt,align=left,font=\scriptsize,text=black!70]
  {one shared head,\\invoked $i=1\dots N$};

% ----- Light backgrounds rectangle grouping the shared per-user stage ----
\begin{scope}[on background layer]
  \node[draw=orange!60,thick,densely dashed,rounded corners=3pt,
        fill=orange!4,inner sep=7mm,
        fit=(FU)(HU)(B1)(BN),
        label={[font=\scriptsize\itshape,text=black!70,anchor=north west]
               north west:Parameter sharing across users}] {};
\end{scope}

\end{tikzpicture}
}
\caption{Two-stage controller. The global head consumes cluster-aggregate features and predicts $(\mu, \chi, \beta_A)$; the per-user head consumes single-user features plus the predicted global block and predicts $\beta_i$. Parameter sharing across users renders the controller dimension independent of $N$.}
\label{fig:ctrl}
\end{figure}

\subsection{Hypothesis Classes}

We compare three deployable model families for each head on the multi-axis trade-off, plus a non-parametric lookup reference (KNN) used only in the single-link comparison:

\subsubsection{Kolmogorov--Arnold Network (KAN)}
KAN \cite{liu2024kan} replaces the fixed pointwise nonlinearities of an MLP with \emph{learnable} univariate functions on every edge, represented as a B-spline plus a residual SiLU basis. The motivation traces back to the Kolmogorov--Arnold representation theorem, which states that any multivariate continuous function on a bounded domain decomposes as a sum of univariate compositions; KAN parameterizes those univariate components as splines and fits them by gradient descent. We use the \texttt{pykan} \cite{liu2024kan2} implementation at width $[d_{\mathrm{in}}, 3, d_{\mathrm{out}}]$, B-spline order 3, and grid size 5. The training pipeline has three stages: (i) standard regression to fit the splines; (ii) pruning of edges whose contribution falls below a threshold; (iii) symbolic candidate fitting from a library of polynomial, $\sin$, $\exp$, and $\tanh$ functions, followed by a joint refit of the network with the chosen symbolic substitution in place. The refit step is critical: skipping it collapses the apparent closed-form to a constant predictor, which was an early pitfall we corrected. The closed-form rule for $\beta$ in Sec.~\ref{sec:results:singlelink} is then read off the trained splines using the methodology of \cite{cranmer2020discovering}.

We state at the outset what the controlled KAN-versus-MLP literature leads us to expect, because our own results land exactly where it predicts. Comparing the two families at matched parameter counts and FLOPs across machine learning, vision, audio and language tasks, \cite{yu2024fairer} reports that MLPs generally outperform KANs \emph{except} on symbolic formula representation, and \cite{pourkamali2024lowdata} finds the parameter overhead of KANs a liability in data-scarce regimes. Our setting is data-scarce by construction, since each training pair costs one call to the decomposed oracle. We therefore do not expect, and do not find, a decisive regression advantage for the KAN over a regularized MLP on the one non-degenerate output; what we do find is the symbolic extraction that \cite{yu2024fairer} identifies as the KAN's remaining advantage, together with a parameter count $43\%$ of the MLP's. Reading our result as a specific instance of that pattern is more defensible than reading it as a general win, and we present it that way.

\subsubsection{Regularized Multi-Layer Perceptron (MLP)}
Two hidden layers of width 32, ReLU, weight decay $10^{-3}$, early stopping on a held-out 20\% split. Adam at $10^{-2}$, 1500 epochs. This is a deliberately strong baseline; an unregularized MLP overfits the small dataset and produces unfair comparisons (we verified this experimentally before fixing the protocol; see Sec.~\ref{sec:results:singlelink}).

\subsubsection{Linear Residual}
A linear least-squares fit of all outputs against all features; serves as the residual base for the KAN/MLP heads.

\subsubsection{$k$-Nearest-Neighbor Lookup (KNN)}
A non-parametric table baseline that predicts each output as the distance-weighted average of its $k{=}5$ nearest neighbors (Euclidean, standardized features) in the training set. It is included in the single-link comparison as a memorization reference: it interpolates the oracle well in regression but, being a lookup table, offers neither parameter economy nor an interpretable rule.

\subsection{Residual Learning}

Direct regression on the cluster-controller dataset proved fragile. Regression MAE remained small (between $10^{-3}$ and $10^{-2}$), yet closed-loop key retention collapsed to zero, because the oracle sat on a sharp 4-D feasibility boundary in $(\mu, \chi, \beta_A, \beta_i)$ space, and multi-directional errors from KAN and MLP pushed predictions across that boundary. The linear baseline, ironically, attained 0.56 key retention by virtue of its smoothing bias toward the cluster mean.

We therefore adopt a \emph{residual learning} strategy. The linear baseline $\mathcal{H}_{\mathrm{lin}}$ is fit first by ordinary least squares; KAN or MLP then regresses on the residual $y - \mathcal{H}_{\mathrm{lin}}(x)$, and the deployed predictor is the sum $\mathcal{H}_{\mathrm{lin}}(x) + \mathcal{H}_{\mathrm{resid}}(x)$. The output layer of the residual network is zero-initialized, so the initial prediction is exactly the linear one; training then carves out a corrective signal only where the linear approximation hurts feasibility. Empirically this lifts feasibility for every controller (Sec.~\ref{sec:results:cluster}). We note that what an earlier draft of this paper called a \emph{boundary collapse} of KAN and MLP was not a property of the feasibility boundary at all but the effect of two software faults, identified and corrected in Sec.~\ref{sec:results:cluster}; the margin recipe is still useful, but it is not rescuing the controllers from a modelling pathology without sacrificing the linear baseline's smoothing bias. The same idea has been used in classical signal processing under the name \emph{boosted regression}, and in physical-layer deep learning under \emph{residual model-based augmentation} \cite{oshea2017introduction,wang2017deep}.

\subsection{$\beta$ Safety-Margin Deployment Recipe}\label{sec:controller:margin}

Even after training, predicted $\hat\beta$ values that fall just below the BSA $m^{\min}$ floor forfeit the entire pair-step. We deploy $\beta_{\mathrm{deploy}} = \hat\beta + \delta$ with a single scalar $\delta^\star$ chosen by 1-D grid sweep on a held-out fold. The curve $\delta\mapsto \mathrm{KeyRetention}$ is unimodal across all controllers, with $\delta^\star\in[0.03, 0.10]$, lifting feasibility by $1.1$--$1.5\times$. The recipe is model-agnostic and adds zero inference cost.

\subsection{Five-Axis Evaluation Protocol}

Each controller is scored on five axes:

\begin{description}
\item[$M_1$:] mean absolute error per output, mean$\pm$std over five seeds.
\item[$M_2$:] closed-loop secret-key retention: ratio of achieved key (controller predictions plugged into the analytic environment) to oracle key (decomposed solver), feasible-gated.
\item[$M_3$:] parameter count: total trainable parameters (heads combined).
\item[$M_4$:] inference latency on a single CPU core, $\mu$s per prediction.
\item[$M_5$:] interpretability: whether a closed-form rule with $R^2 \ge 0.95$ can be extracted (boolean per output).
\end{description}

We deliberately report all five rather than collapsing to a single ``score,'' because the controllers occupy a Pareto front; the right model for deployment depends on the constraint mix. Interpretability is decisive for safety-critical satellite firmware, while raw key retention dominates for an airborne ground link, and parameter count matters for an embedded payload.
